% ====================================================================
% FF2-4 (ch1 §15 식 eq:U1T2 아래 배치 제안) — T1 중심 U_1(T) 의 전자항 T² 곡률 식별 신호
%   좌표 = eq:U1T2 를 python 으로 점별 수치 평가:
%   ΔS_0=+80 J/(mol K)(대표 스케일, tier B — §12 verifybox), a_e=−(π²/3)R(k_B/e_V)(g_max/Δx)·¼
%   = −0.15321 J/(mol K²) (게이트 중심 — eq:dSegate·eq:dSemolar 닫힌식; a_e·T_ref=−45.7≈−46 J/(mol K)),
%   T_ref=298.15 K. (b) 기울기 공간: ∂U_1/∂T = (ΔS_0+a_e T)/F — T 에 선형(식별 신호).
%   사용 라이브러리: ch1 preamble 범위 내(arrows.meta).
% ====================================================================
\begin{figure}[t]
\centering
\begin{tikzpicture}
% ============ (a) U_1(T) - U_1(T_ref) ============
\begin{scope}[x=0.055cm,y=0.05cm,shift={(-268.15,0)}]
\draw[-{Latex[length=1.6mm]}] (266,-33) -- (266,50)
  node[anchor=south west,xshift=-4mm,font=\scriptsize] {$U_1(T)-U_1(T_\mathrm{ref})$ [mV]};
\draw[-{Latex[length=1.6mm]}] (266,-33) -- (356,-33)
  node[anchor=north west,xshift=-3mm,font=\scriptsize] {$T$ [K]};
\foreach \xx in {270,290,310,330,350} {\draw (\xx,-31.6) -- (\xx,-34.4) node[below,font=\scriptsize] {\xx};}
\foreach \yy in {-30,-15,15,30,45} {\draw (266.6,\yy) -- (265.4,\yy) node[left,font=\scriptsize] {\yy};}
\draw[densely dashed,black!45] (266,0) -- (352,0);
\draw[densely dotted,black!55] (298.15,-33) -- (298.15,46);
\node[font=\tiny,black!55,anchor=south] at (298.15,46) {$T_\mathrm{ref}$};
% baseline linear (no electronic term) — straight line
\draw[thick,densely dotted] (268.15,-24.87) -- (348.15,41.46);
% full T^2 curvature (eq:U1T2)
\draw[very thick] plot[smooth] coordinates {(268.15,-11.39) (272.15,-9.78) (276.15,-8.21) (280.15,-6.66) (284.15,-5.14) (288.15,-3.64) (292.15,-2.16) (296.15,-0.71) (300.15,0.71) (304.15,2.11) (308.15,3.48) (312.15,4.82) (316.15,6.15) (320.15,7.44) (324.15,8.71) (328.15,9.96) (332.15,11.18) (336.15,12.37) (340.15,13.54) (344.15,14.68) (348.15,15.80)};
% right-end stacked labels + identification gap arrow
\draw[{Latex[length=1.3mm]}-{Latex[length=1.3mm]}] (349.4,41.46) -- (349.4,15.80);
\node[font=\tiny,anchor=west] at (350.4,41.46) {기저 선형 $0.83$ mV/K};
\node[font=\tiny,anchor=west] at (350.4,28.6) {$-25.7$ mV};
\node[font=\tiny,anchor=west,align=left] at (350.4,15.80) {적분형 $T^2$\\(식~\eqref{eq:U1T2})};
\node[font=\scriptsize,anchor=north,align=center] at (306,-13) {고온 외삽이 선형보다\\\emph{낮아짐} $=$ 전자항 식별 신호};
\node[font=\scriptsize] at (310,-42) {(a)};
\end{scope}
% ============ (b) slope space: dU_1/dT(T) ============
\begin{scope}[xshift=8.3cm,x=0.055cm,y=4.6cm,shift={(-268.15,-0.16)}]
\draw[-{Latex[length=1.6mm]}] (266,0.16) -- (266,0.98)
  node[anchor=south west,xshift=-4mm,font=\scriptsize] {$\partial U_1/\partial T$ [mV/K]};
\draw[-{Latex[length=1.6mm]}] (266,0.16) -- (356,0.16)
  node[anchor=north west,xshift=-3mm,font=\scriptsize] {$T$ [K]};
\foreach \xx in {270,290,310,330,350} {\draw (\xx,0.168) -- (\xx,0.152) node[below,font=\scriptsize] {\xx};}
\foreach \yy in {0.2,0.4,0.6,0.8} {\draw (266.6,\yy) -- (265.4,\yy) node[left,font=\scriptsize] {\yy};}
\draw[densely dotted,black!55] (298.15,0.16) -- (298.15,0.90);
\node[font=\tiny,black!55,anchor=south] at (298.15,0.90) {$T_\mathrm{ref}$};
% baseline (no electronic): constant
\draw[thick,densely dotted] (268.15,0.829) -- (348.15,0.829);
\node[font=\tiny,anchor=south] at (308,0.836) {기저(전자항 0): $\Delta S_0/F=0.829$ (상수)};
% frozen approximation: constant at T_ref value
\draw[thick,densely dashed] (268.15,0.3557) -- (352,0.3557);
\node[font=\tiny,anchor=west] at (330,0.322) {동결 근사 $0.356$ (상수)};
\draw[-{Latex[length=1.1mm]},black!60] (335,0.335) -- (337,0.353);
% full: linear in T
\draw[very thick] (268.15,0.4033) -- (348.15,0.2763);
\fill (298.15,0.3557) circle (1.2pt);
\node[font=\tiny,anchor=south west,align=left] at (300.5,0.545)
  {적분형: $(\Delta S_0{+}a_eT)/F$\\--- $T$ 에 \emph{선형}(기울기 $a_e/F$)};
\draw[-{Latex[length=1.2mm]},black!60] (313,0.52) -- (321.5,0.324);
\node[font=\tiny,anchor=west] at (269,0.452) {$0.403$};
\node[font=\tiny,anchor=north] at (344,0.262) {$0.276$};
\node[font=\tiny,black!60,anchor=north east] at (297,0.30) {$T_\mathrm{ref}$ 에서 동결 근사와 접함};
\node[font=\scriptsize] at (310,0.085) {(b)};
\end{scope}
\end{tikzpicture}
\caption{T1 전이 중심의 온도 함수형이 갖는 전자항 식별 신호 --- 좌표는 식~\eqref{eq:U1T2} 그대로의
수치 평가다($\Delta S_0{=}{+}80$ J/(mol\,K) 대표 스케일[tier B --- \S\ref{sec:lco-center}], 게이트
중심의 몰당 닫힌식(식~\eqref{eq:dSegate}$\cdot$\eqref{eq:dSemolar})에서 $a_e{=}{-}0.153$
J/(mol\,K$^2$), $T_\mathrm{ref}{=}298.15$ K). (a) $U_1(T)$: 전자항이 없으면 기저 선형(점선,
$0.83$ mV/K), 전자항($\propto T$)을 적분하면 $T^2$ 곡률(실선) --- $a_e<0$ 라 고온 외삽이 선형
외삽보다 \emph{낮아지고}($348$ K 에서 $-25.7$ mV) 이것이 다온도 피팅의 식별 신호다. (b) 같은 내용의
기울기 공간: $\partial U_1/\partial T=(\Delta S_0+a_eT)/F$ 가 $T$ 에 \emph{선형}으로 감소한다
($0.403\to0.276$ mV/K). 현행 동결 근사(파선 --- $\Delta S_e$ 를 $T_\mathrm{ref}$ 상수로 두는
\S\ref{sec:lco-code} 의 단일-기준 근사)는 적분형의 $T_\mathrm{ref}$ 접선이라 (a)에서는 실선과
$(a_e/2F)(T-T_\mathrm{ref})^2$($348$ K 에서 $-2.0$ mV)만큼만 이차로 벌어지고, (b)의 기울기
공간에서 상수 대 선형으로 명확히 갈린다.}
\label{fig:cand-ff2-4}
\end{figure}
